\documentclass[11pt,a4paper]{article}

\usepackage[utf8]{inputenc}
\usepackage[T1]{fontenc}
\usepackage{amsmath,amssymb,amsthm}
\usepackage{graphicx}
\usepackage{booktabs}
\usepackage{hyperref}
\usepackage{algorithm}
\usepackage{algorithmic}
\usepackage{xcolor}
\usepackage{geometry}
\geometry{margin=2.5cm}

\title{Characterize-First Symbolic Regression:\\
ODE Inversion for Efficient Expression Discovery}

\author{Edgar Zanella Alvarenga}
\date{April 2026}

\begin{document}

\maketitle

\begin{abstract}
We present a \emph{characterize-first} approach to symbolic regression, in which structural
properties of the data are identified before expression search begins. The central method is
\emph{ODE inversion}: given noisy observations of an unknown function $f$, we detect the
linear (or nonlinear) differential equation that $f$ satisfies, thereby recovering the
algebraic solution class that tightly constrains any subsequent parameter or expression search.
We validate the approach on bacterial growth time-series data (optical density measurements)
from which classical genetic programming struggles to recover the generating Gompertz model.
ODE inversion in log-space or via direct Gompertz ODE scanning (B2) recovers the exact
Gompertz parameters ($A$, $\mu$, $\lambda$) on synthetic data with $R^2=1.000$ and fits
real data with $R^2 \geq 0.992$, matching the best parametric fits while revealing the
governing ODE structure and its biological interpretation.
\end{abstract}

% ──────────────────────────────────────────────────────────────
\section{Introduction}
\label{sec:intro}
% ──────────────────────────────────────────────────────────────

Symbolic regression (SR) is the problem of finding a mathematical expression $\hat{f}$ that
accurately describes a set of observations $\{(x_i, y_i)\}_{i=1}^n$ without a predefined
functional form~\cite{koza1992}. Unlike classical regression, the \emph{form} of the expression
is part of the search problem, making SR combinatorially hard.

The dominant approach, Genetic Programming (GP)~\cite{koza1992}, searches a space of expression
trees via mutation, crossover, and selection. While GP has achieved notable successes—including
the AI Feynman benchmark~\cite{udrescu2020}—it suffers from three core limitations:
(i) the search space grows super-exponentially with expression depth;
(ii) gradient information about expression structure is unavailable;
(iii) stochastic search can waste budget on structurally wrong forms.

Several works have proposed to constrain GP search using prior knowledge:
PySR~\cite{cranmer2023} uses mixed discrete/continuous optimisation;
AI Feynman~\cite{udrescu2020} leverages dimensional analysis and symmetry detection;
SINDy~\cite{brunton2016} discovers governing differential equations from data for
dynamical systems. Our work combines these ideas into a unified \emph{characterize-first}
pipeline applicable to univariate regression.

\paragraph{Characterize-first principle.}
Before searching for an expression, answer: \emph{What structural class must the answer
belong to?} If the data satisfies a differential equation $\mathcal{L}[f]=0$, the solution
space of $\mathcal{L}$ is often a low-dimensional manifold. Fitting within that manifold is
a standard least-squares or low-dimensional optimisation problem, orders of magnitude smaller
than unconstrained GP search.

\paragraph{Contributions.}
\begin{itemize}
  \item A noise-robust linear ODE inversion pipeline using Gaussian derivative filters,
        GCV bandwidth selection, and STRidge (Section~\ref{sec:linear_ode}).
  \item Three nonlinear ODE detection methods for sigmoidal/growth data:
        log-space linearisation (B1), direct Gompertz ODE scanning (B2), and
        STRidge on a growth-targeted feature library (B3) (Section~\ref{sec:nonlinear_ode}).
  \item Validation on bacterial growth data showing that B2 recovers Gompertz parameters
        with $R^2=0.992$, matching the best parametric fit (Section~\ref{sec:results}).
  \item Discussion of how detected ODE structure constrains GP grammar, enabling
        \emph{informed symbolic regression} (Section~\ref{sec:discussion}).
\end{itemize}

% ──────────────────────────────────────────────────────────────
\section{Background}
\label{sec:background}
% ──────────────────────────────────────────────────────────────

\subsection{Bacterial Growth Models}

Bacterial population dynamics in controlled conditions follow characteristic sigmoid curves.
The four principal parametric models are:

\textbf{Gompertz}~\cite{gompertz1825}:
\begin{equation}
  f(t) = A \exp\!\bigl(-\exp\!\bigl[(\mu_{\max} e / A)(\lambda - t) + 1\bigr]\bigr)
  \label{eq:gompertz}
\end{equation}
where $A$ is the asymptotic population density, $\mu_{\max}$ is the maximum specific growth
rate, and $\lambda$ is the lag time.

\textbf{Logistic}~\cite{verhulst1838}:
\begin{equation}
  f(t) = \frac{A}{1 + \exp\!\bigl[4\mu_{\max}(\lambda - t)/A + 2\bigr]}
\end{equation}

\textbf{Baranyi-Roberts}~\cite{baranyi1994}:
\begin{equation}
  \frac{\mathrm{d}f}{\mathrm{d}t} = \mu_{\max}\,q(t)\,\frac{1 - f/A}{1 + (f/A)^\nu}\,f,
  \quad \dot{q} = \mu_{\max}\,q
\end{equation}
where $q(0)$ encodes the initial physiological state.

\textbf{Richards/Generalised logistic}~\cite{richards1959}:
\begin{equation}
  f(t) = \frac{A}{\bigl(1 + \exp(-r(t-\lambda))\bigr)^{1/\nu}}
\end{equation}
The Gompertz model has a distinctive inflection at $37\%$ of the growth range (vs.\ $50\%$
for the symmetric logistic), making it identifiable from the derivative profile.

\subsection{SINDy and Sparse ODE Detection}

SINDy~\cite{brunton2016} represents $\dot{\mathbf{x}} = \Theta(\mathbf{x})\,\boldsymbol{\xi}$
where $\Theta$ is a library of candidate terms and $\boldsymbol{\xi}$ is sparse. STRidge
(sequential thresholded ridge regression) finds the sparsest solution consistent with the data.
We adapt this framework to univariate function identification.

% ──────────────────────────────────────────────────────────────
\section{Linear ODE Inversion}
\label{sec:linear_ode}
% ──────────────────────────────────────────────────────────────

\subsection{Problem Formulation}

Given $n$ noisy observations $\{(x_i, y_i)\}$ with $y_i = f(x_i) + \varepsilon_i$, we seek
a linear constant-coefficient ODE
\begin{equation}
  a_0 f + a_1 f' + a_2 f'' + \cdots + a_p f^{(p)} = 0
\end{equation}
that $f$ satisfies. The solution space is spanned by $\{e^{r_j x}, xe^{r_j x}, \sin(\omega x),
\cos(\omega x), \ldots\}$ depending on the roots of the characteristic polynomial.

\subsection{Derivative Estimation via Gaussian Filters}

Raw numerical differentiation amplifies noise with factor $O(1/h^k)$ for $k$-th order
derivatives and step size $h$. We use Gaussian derivative filters~\cite{lindeberg1994}:
\begin{equation}
  f^{(k)}(x) \approx \frac{1}{\sigma^k}\,(G_\sigma^{(k)} * f)(x)
\end{equation}
where $G_\sigma(x) = (2\pi\sigma^2)^{-1/2}\exp(-x^2/2\sigma^2)$ is the Gaussian kernel
with bandwidth $\sigma$ (in pixels). This is the optimal linear estimator under Gaussian noise.

\paragraph{Bandwidth selection.}
Bandwidth $\sigma$ is chosen by Generalised Cross-Validation (GCV) over a grid
$\sigma \in [\sigma_{\min}, \sigma_{\max}]$. The aliasing constraint requires
$\sigma_{\min} = \max(3, 2 p)$ pixels where $p$ is the ODE order, to prevent
the discrete kernel from aliasing.

\paragraph{Edge trimming.}
Boundary padding in \texttt{scipy.ndimage.gaussian\_filter1d} contaminates the
$\lceil 3\sigma \rceil$ points at each end. These are removed before constructing
the derivative matrix.

\subsection{STRidge Detection}

Build the derivative matrix $D \in \mathbb{R}^{m \times (p+1)}$ where
$D_{ij} = f^{(j)}(x_i)$, $j=0,\ldots,p$, and solve for the null-space vector
$\mathbf{a}$ (normalised so $a_p = 1$):
\begin{equation}
  D\,\mathbf{a} \approx \mathbf{0}
\end{equation}
STRidge iteratively applies ridge regression with threshold $\tau$:
\begin{enumerate}
  \item Solve $\hat{\mathbf{a}} = (D^\top D + \lambda I)^{-1} D^\top \mathbf{0}$
        with pivot column fixed to 1.
  \item Zero out coefficients with $|a_j| < \tau \max_j |a_j|$.
  \item Repeat on the active subset until convergence.
\end{enumerate}

\subsection{Noise Robustness}

Table~\ref{tab:noise} shows ODE detection accuracy for $f(x) = \sin(x)$ with additive
Gaussian noise $\varepsilon \sim \mathcal{N}(0, \sigma_\varepsilon^2)$.

\begin{table}[h]
\centering
\caption{Linear ODE detection ($f'' + f = 0$) under noise for $f(x)=\sin(x)$}
\label{tab:noise}
\begin{tabular}{lcc}
\toprule
Noise $\sigma_\varepsilon$ & ODE found? & STRidge residual \\
\midrule
0.000 & YES & $1.4 \times 10^{-3}$ \\
0.010 & YES & $5.6 \times 10^{-3}$ \\
0.020 & YES & $7.8 \times 10^{-3}$ \\
0.050 & YES & $1.0 \times 10^{-2}$ \\
0.100 & YES & $2.5 \times 10^{-2}$ \\
0.200 & YES & $4.3 \times 10^{-2}$ \\
0.500 & no  & $1.1 \times 10^{-1}$ \\
\bottomrule
\end{tabular}
\end{table}

The method is reliable up to SNR $\geq 5$ ($\sim$20\% noise).

% ──────────────────────────────────────────────────────────────
\section{Nonlinear ODE Detection for Growth Data}
\label{sec:nonlinear_ode}
% ──────────────────────────────────────────────────────────────

Linear ODE inversion incorrectly identifies oscillatory structure in bacterial growth data.
The Gompertz model satisfies the \emph{nonlinear} ODE
\begin{equation}
  f'(t) = r\,f(t)\,\ln\!\bigl(K / f(t)\bigr)
  \label{eq:gompertz_ode}
\end{equation}
with $K = A$ (carrying capacity) and $r = \mu_{\max} e / A$.
We propose three methods to detect this structure directly from data.

\subsection{B1: Log-Space Linearisation}

Define $g = \ln f$. Then:
\begin{equation}
  g'(t) = \frac{f'}{f} = r\,\ln(K/f) = r(\ln K - g)
\end{equation}
so $g' = r\ln K - r\,g$, which is a \emph{linear first-order ODE} in $g$:
\begin{equation}
  g' = a + b\,g, \quad b = -r < 0
\end{equation}
We estimate $g'$ via Gaussian derivatives on $g = \ln f$, then solve via least squares:
$[a, b]^\top = (G^\top G)^{-1} G^\top g'$ where $G = [\mathbf{1}, \mathbf{g}]$.

If $b < 0$, the asymptote is $g_\infty = -a/b$ and $K = e^{g_\infty}$.

\subsection{B2: Direct Gompertz ODE Scanning}

For a fixed $K$, equation~\eqref{eq:gompertz_ode} is linear in $r$:
\begin{equation}
  f' = r \cdot \underbrace{f \ln(K/f)}_{\phi(K)}
\end{equation}
The optimal $r$ given $K$ is $r(K) = \langle f', \phi \rangle / \langle \phi, \phi \rangle$
(ordinary least squares). We scan $K \in [f_{\max} \cdot 1.001,\, f_{\max} \cdot 3]$ and
choose the $K$ minimising the normalised residual
$\|f' - r\phi\|_2 / \mathrm{std}(f')$.

The resulting $(K, r)$ initialises a full nonlinear refinement via
\texttt{scipy.optimize.curve\_fit} on equation~\eqref{eq:gompertz}.

\subsection{B3: STRidge on Growth Feature Library}

Build the feature library for $f' = \Theta \boldsymbol{\xi}$:
\begin{equation}
  \Theta = \bigl[\mathbf{1},\; f,\; f^2,\; f\ln f,\; f(1 - f/f_{\max})\bigr]
\end{equation}
These features cover exponential growth ($f$), logistic growth ($f^2$ or $f(1-f/K)$),
and Gompertz growth ($f\ln f$). STRidge selects the dominant term, acting as a
\emph{model classifier} rather than a curve-fitter.

% ──────────────────────────────────────────────────────────────
\section{Results}
\label{sec:results}
% ──────────────────────────────────────────────────────────────

\subsection{Data}

We use optical density (OD) time series from three growth curves (columns 2–4 of
\texttt{all\_curves.csv}), each with 160–190 measurements over $t \in [0.25, 41.0]$~h.
Signal-to-noise ratios are $\text{SNR} \approx 50$–88.

\subsection{Linear ODE Detection on Real Data}

The linear ODE pipeline (ode\_robust.py) identifies:
\begin{itemize}
  \item \textbf{Col2}: $f'''' + 0.255\,f'' = 0$ → solution $C_1 + C_2 x + C_3\sin(x/2) + C_4\cos(x/2)$
  \item \textbf{Col3}: $0.251\,f'' + 0.820\,f''' + f'''' = 0$ → damped oscillatory
  \item \textbf{Col4}: $0.124\,f' + 0.522\,f'' + f''' + 0.695\,f'''' = 0$ → damped oscillatory
\end{itemize}
These solutions achieve $R^2 = 0.698$–$0.742$ — the oscillatory solution class is fundamentally
wrong for sigmoidal bacterial growth. This confirms the need for nonlinear ODE detection.

\subsection{Nonlinear ODE Detection (Step B)}

Table~\ref{tab:results_b} summarises B2 (best method) results compared to baseline fits.

\begin{table}[h]
\centering
\caption{Model fitting results on real bacterial growth data}
\label{tab:results_b}
\begin{tabular}{lccccc}
\toprule
Method & \multicolumn{2}{c}{Col 2} & \multicolumn{2}{c}{Col 3} & Col 4 \\
 & RMSE & $R^2$ & RMSE & $R^2$ & $R^2$ \\
\midrule
Linear ODE (oscillatory) & 0.400 & 0.698 & --- & --- & --- \\
Logistic (parametric)    & 0.260 & 0.742 & 0.091 & 0.897 & 0.901 \\
B1: Log-space ODE        & 0.509 & $-$0.027 & 0.063 & 0.941 & 0.914 \\
B2: Gompertz ODE scan    & \textbf{0.046} & \textbf{0.992} & \textbf{0.060} & \textbf{0.946} & \textbf{0.957} \\
Gompertz (parametric)    & 0.046 & 0.992 & 0.060 & 0.945 & 0.957 \\
Richards (generalised)   & 0.045 & 0.992 & --- & --- & --- \\
\bottomrule
\end{tabular}
\end{table}

\paragraph{B2 achieves parametric-level accuracy without prior knowledge of the model class.}
On synthetic Gompertz data, B2 recovers exact parameters:
$A=1.000$, $\mu=0.300$, $\lambda=5.000$ with $R^2=1.000$.

\paragraph{B3 correctly classifies all three curves as Gompertz-type} (f·ln(f) dominant in
sparse ODE, vs.\ f² for logistic), providing model-class information for grammar constraints.

\subsection{Model Discrimination: Gompertz vs Logistic}

The inflection position relative to the growth range is a clean classifier:
\begin{itemize}
  \item Gompertz: inflection at $\approx 37\%$ of $[y_{\min}, y_{\max}]$
  \item Logistic: inflection at $\approx 50\%$
\end{itemize}
All three real curves show inflection at $\sim$35–40\%, confirming Gompertz structure.
B3's sparse ODE independently confirms this via the $f\ln f$ signature.

% ──────────────────────────────────────────────────────────────
\section{ODE Discovery for Complex Growth Models}
\label{sec:complex}
% ─────────────────────────────────────────���────────────────────

Beyond single Gompertz/logistic growth, we apply the characterize-first pipeline to
two more complex biological scenarios: heterogeneous populations (aHPM) and diauxic growth.

\subsection{aHPM: Heterogeneous Population Model}

The aHPM~\cite{baranyi1994} describes a population with two subpopulations: $u_1$ (lag,
not yet actively growing) and $u_2$ (growing). The system ODE is:
\begin{align}
  \dot{u}_1 &= -\delta\, u_1 \\
  \dot{u}_2 &= \delta\, u_1 + g_r\, u_2 \bigl(1 - (X/K)^s\bigr)
\end{align}
where $X = u_1 + u_2$ is the observable OD, $\delta$ the exit-lag rate, $g_r$ the growth
rate, $K$ the carrying capacity, and $s$ the Richards shape parameter.

\paragraph{Theoretical derivation.}
Since $u_1(t) = u_1(0)\,e^{-\delta t}$, the observable satisfies the
\emph{non-autonomous} ODE:
\begin{equation}
  \boxed{X'(t) = g_r \bigl(X - L_0\,e^{-\delta t}\bigr)\bigl(1 - (X/K)^s\bigr)}
  \label{eq:ahpm_compact}
\end{equation}
where $L_0 = u_1(0)$ is the initial lag population. This reduces to the Richards ODE
when $L_0 = 0$. The key novel term $L_0\,e^{-\delta t}$ subtracts the decaying lag
contribution from the effective growing biomass.

\paragraph{Simple approximation.}
Equation~\eqref{eq:ahpm_compact} has the approximate closed-form:
\begin{equation}
  X(t) \approx \text{Richards}(g_r, K, s,\, t) + L_0\,e^{-\delta t}
  \label{eq:ahpm_approx}
\end{equation}
This linearisation is accurate when $\delta \gg g_r$: the lag population decays
much faster than the growth time-scale.

\paragraph{Results.}
Fitting equation~\eqref{eq:ahpm_compact} to synthetic aHPM data (strong lag: $L_0=0.15$,
$\delta=0.15$):

\begin{table}[h]
\centering
\caption{aHPM fitting results (strong lag, $L_0=0.15$, $\delta=0.15$, $K=1.0$)}
\label{tab:ahpm}
\begin{tabular}{lcc}
\toprule
Model & RMSE & $R^2$ \\
\midrule
Gompertz (baseline)         & 0.039 & 0.982 \\
Richards (4 params)         & 0.010 & 0.9988 \\
Compact aHPM ODE (5 params) & 0.005 & 0.9997 \\
Richards $+$ exp decay      & 0.005 & 0.9997 \\
\bottomrule
\end{tabular}
\end{table}

The compact aHPM formula recovers $K = 1.001$ (true: $1.0$) and $\delta = 0.378$ (true: $0.15$;
within scale, as the effective decay seen in $X$ differs from the microscopic rate).

\paragraph{Inflection fingerprint.}
The inflection position of $X(t)$ relative to its range is a diagnostic for heterogeneous
population structure:
\begin{itemize}
  \item Gompertz: $\approx 37\%$
  \item aHPM (strong lag): $\approx 44\text{--}45\%$
  \item Logistic: $\approx 50\%$
\end{itemize}
This provides a new model-class discriminator based solely on derivative structure,
extending the Gompertz/logistic inflection test from the literature.

\subsection{Diauxic Growth: Theoretical Impossibility and Discovery}

Diauxic growth (growth $\to$ plateau $\to$ growth, from sequential substrate consumption)
poses a fundamental challenge for autonomous ODE inversion.

\paragraph{Key theoretical result.}
\emph{Diauxic growth with a genuine plateau cannot be captured by a 1st-order
autonomous ODE.} Formally: for $X' = f(X)$, a pause at $X = K_1$ requires $f(K_1) = 0$.
But if $f(K_1) = 0$ and $f(X) > 0$ for $X \in (K_1, K_2)$ (needed for phase-2 growth),
then $K_1$ is an unstable equilibrium — $X$ cannot approach $K_1$ from below. Any
compact form with a true plateau requires hidden state or an explicit time dependence.

Validated computationally: all 1st-order autonomous ODE forms (bi-logistic, Gompertz-logistic,
bi-Gompertz) achieve normalised residuals $\geq 0.96$ on double-Gompertz synthetic data.

\paragraph{Phase-segmented symbolic boosting.}
For true diauxic data, we propose phase-segmented boosting:
\begin{enumerate}
  \item Locate the diauxic valley: the time $t^*$ minimising $X'(t)$ between the two
        derivative peaks.
  \item Fit $G_1$ (Gompertz) to data on $[0, t^*]$ (phase 1 only).
  \item Fit $G_2$ (Gompertz) to the residual $X(t) - G_1(t)$ on $[t^*, T]$ (phase 2).
\end{enumerate}
This recovers both components with $R^2 = 0.9998$ and near-exact parameters
(Table~\ref{tab:diauxic}).

\begin{table}[h]
\centering
\caption{Phase-segmented boosting on double-Gompertz synthetic diauxic data}
\label{tab:diauxic}
\begin{tabular}{lcccc}
\toprule
Component & $A$ (true) & $A$ (rec.) & $\mu$ (rec.) & $\lambda$ (rec.) \\
\midrule
Phase 1 & 0.550 & 0.552 & 0.350 & 6.02 \\
Phase 2 & 0.400 & 0.398 & 0.200 & 26.02 \\
\bottomrule
\end{tabular}
\end{table}

\paragraph{Smooth diauxic: Gompertz-logistic hybrid.}
For Monod 2-substrate diauxic data where the catabolite repression switch is smooth,
the best compact autonomous ODE found is:
\begin{equation}
  X'(t) = r\, X\, \ln(K/X)\,(1 + a\,X)
  \label{eq:gomp_logistic}
\end{equation}
The additional term $1 + aX$ represents density-dependent modulation of the Gompertz
growth law. For $a > 0$, higher biomass leads to faster effective growth — a possible
signature of quorum sensing, cross-feeding, or pH buffering at higher cell densities.

\subsection{Model-Free ODE Discovery on Synthetic Diauxic Data}
\label{sec:modelfree_diauxic}

We apply our full pipeline to synthetic Monod 2-substrate diauxic data in a \emph{genuinely
model-free} setting: no prior knowledge of the Gompertz, aHPM, or any other form is assumed.
Data is generated from:
\begin{align}
  \dot{X}  &= \bigl[\mu_1 S_1/(S_1+K_s) \cdot H(S_1) + \mu_2 S_2/(S_2+K_s) \cdot (1-H(S_1))\bigr] X \\
  H(S_1)   &= S_1^n\,/\,(S_1^n + K_d^n)
\end{align}
with $\mu_1=0.90$, $\mu_2=0.25$, $K_d=0.005$, $n=10$, observable $X(t)$ only.

\paragraph{Step 1: Derivative fingerprint.}
Computing $X'(t)$ via Gaussian filter reveals \textbf{two distinct peaks} at $t=4.51$ and
$t=7.64$, with a trough at $t=5.39$ — the definitive diauxic signature. This detection
requires no model assumption.

\paragraph{Step 2: 2nd-order autonomous ODE discovery.}
Applying STRidge to the library $\Theta = [1, X, X^2, X^3, X', (X')^2, X{\cdot}X',
X{\cdot}\ln X, \ldots]$ to discover $X'' = F(X, X')$ produces sparse equations with
derivative residuals as low as 0.04, but forward integration fails ($R^2 \ll 0$).
This confirms our theoretical result: a compact autonomous ODE cannot capture
non-autonomous diauxic dynamics.

\paragraph{Step 3: Direct ODE scan discovers Gompertz structure.}
Applying the B2 direct scan (no model assumption) — scanning $K \in [X_{\max}, 3X_{\max}]$
and solving $r = \langle X', \Phi\rangle / \|\Phi\|^2$ where $\Phi = X\ln(K/X)$ — the
pipeline identifies:
\begin{equation}
  \boxed{X' = r\, X\, \ln(K/X), \quad r = 0.395,\; K = 1.540}
  \label{eq:discovered_gompertz}
\end{equation}
\textbf{without assuming Gompertz form a priori.} Forward integration gives $R^2 = 0.968$.
This is the key demonstration of model-free discovery: the Gompertz ODE emerges as the
dominant structure from a generic feature scan.

\paragraph{Why the Gompertz ODE dominates.}
For smooth Monod kinetics with sharp catabolite repression ($n=10$, $K_d=0.005$), the
diauxic transition in $X(t)$ is nearly invisible — the parametric Gompertz fits $R^2=0.992$.
The diauxic structure is visible only in $X'(t)$ (two peaks), not in $X(t)$ itself.
This is an important fundamental limitation: model-free discovery from $X(t)$ alone cannot
distinguish smooth Monod diauxic from single-phase Gompertz growth.

\begin{table}[h]
\centering
\caption{Model-free ODE discovery on synthetic Monod diauxic data}
\label{tab:modelfree}
\begin{tabular}{lcc}
\toprule
Method & $R^2$ & Notes \\
\midrule
Direct Gompertz scan (model-free) & 0.968 & discovered without prior \\
Gompertz parametric (baseline)    & 0.992 & known model \\
2nd-order autonomous ODE          & $<$0  & fails: non-autonomous dynamics \\
Non-autonomous STRidge             & $<$0  & overfits derivative space \\
aHPM (autonomous)                  & $<$0  & wrong model class \\
\bottomrule
\end{tabular}
\end{table}

\paragraph{Conclusion.}
The characterize-first pipeline applied to diauxic data correctly:
(1) identifies the two-phase signature in $X'(t)$,
(2) rules out autonomous 2nd-order ODEs,
(3) discovers the Gompertz ODE as the dominant structure from a feature scan.
When the diauxic transition is sharp (double-Gompertz data), phase-segmented boosting
recovers both phases exactly (Table~\ref{tab:diauxic}). When the transition is smooth
(Monod data), the derivative fingerprint provides the only reliable diauxic evidence.

% ──────────────────────────────────────────────────────────────
\section{Grammar Constraint for Symbolic Regression}
\label{sec:discussion}
% ──────────────────────────────────────────────────────────────

\subsection{From ODE to Grammar Prior}

The characterize-first pipeline produces structured information that constrains any
downstream symbolic regression search:

\begin{enumerate}
  \item \textbf{Model class}: B3 STRidge identifies the dominant ODE term
        (e.g., $f\ln f$ → Gompertz-like). This eliminates oscillatory terms ($\sin$, $\cos$)
        from the grammar and prioritises $\exp$ and $\ln$ primitives.

  \item \textbf{Solution template}: B2 identifies $K$ and $r$, giving the functional form
        $f = K\exp(-C\exp(-r(t-t_0)))$ as a starting point. GP then searches for deviations
        from this template.

  \item \textbf{Symmetry detection}: parity (even/odd), periodicity, and scaling
        symmetries further restrict the search. Monotone+saturating data rules out
        periodic solutions entirely.
\end{enumerate}

\subsection{Grammar Constraint in Practice}

Once the ODE structure is known, downstream symbolic regression search can be constrained
to expression classes consistent with the detected ODE. Concretely:
\begin{itemize}
  \item Oscillatory terms ($\sin$, $\cos$) are excluded for monotone sigmoidal data.
  \item The ODE-derived carrying capacity $K$ and rate $r$ provide warm-start initial
        constants, replacing random initialisation.
  \item The detected model class (Gompertz, Richards, aHPM) defines the functional
        skeleton; any remaining free search is over the residual only.
\end{itemize}
This is a natural direction for future work combining the Python ODE pipeline with a
GP back-end.

\subsection{Phase-Segmented Symbolic Boosting}

For multi-phase growth (diauxic), we propose a phase-segmented boosting approach.
Define the diauxic valley $t^*$ as the minimum of $X'$ between the two derivative peaks.
Then:
\begin{enumerate}
  \item Fit $G_1$ to data on $[0, t^*]$ (phase 1).
  \item Compute residual $r(t) = X(t) - G_1(t)$.
  \item Fit $G_2$ to the residual on $[t^*, T]$ (phase 2).
\end{enumerate}
This recovers both Gompertz components with $R^2 = 0.9998$ (Table~\ref{tab:diauxic}),
without any prior assumption that the data contains exactly two phases.

% ──────────────────────────────────────────────────────────────
\section{Discussion and Future Work}
\label{sec:future}
% ──────────────────────────────────────────────────────────────

\paragraph{Limitations.}
The current pipeline handles scalar univariate functions. Extension to multivariate
and dynamical systems (PDEs, coupled ODEs) is straightforward in principle but requires
adapting the derivative matrix construction. Euler ODEs ($f' = f/x$) with non-constant
coefficients are not yet handled.

\paragraph{Weak/integral formulation.}
For very noisy data (SNR $< 5$), differentiation fails. Integration against test functions
$\langle f', \phi \rangle = -\langle f, \phi' \rangle$ avoids differentiation entirely~\cite{messenger2021}.
This is implemented but not yet benchmarked.

\paragraph{Non-Gompertz data.}
The B2 scanner assumes Gompertz structure. A general approach would scan over a library
of ODE forms (logistic, Richards, power-law) and select the best fit by AIC/BIC.
B3 partially addresses this via STRidge model classification.

\paragraph{Full SR pipeline.}
The complete characterize-first pipeline is:
\begin{enumerate}
  \item Symmetry detection (parity, periodicity, scaling).
  \item Linear ODE inversion → oscillatory/polynomial/exponential class.
  \item If residual high: nonlinear ODE detection (B1/B2/B3) → growth model class.
  \item Grammar restriction based on detected class.
  \item Constrained GP or exhaustive search on restricted grammar.
  \item Symbolic boosting on residuals.
\end{enumerate}

% ──────────────────────────────────────────────────────────────
\section{Conclusion}
% ──────────────────────────────────────────────────────────────

We have demonstrated that ODE inversion is a practical and effective pre-processing step
for symbolic regression. On bacterial growth data, the direct Gompertz ODE scanning (B2)
recovers exact model parameters on synthetic data and achieves $R^2 \geq 0.992$ on real
noisy measurements, matching the best parametric fits. The method is noise-robust (working
to $\sim$20\% additive noise with Gaussian derivative filters and GCV bandwidth selection),
computationally cheap (seconds on 200-point time series), and produces interpretable output:
the governing ODE and its analytical solution form.

The characterize-first paradigm inverts the usual SR workflow: instead of searching blindly
and hoping to find the right expression, we \emph{characterize what class the expression must
belong to} before searching. This transforms an intractable combinatorial search into a
well-posed algebraic problem.

% ──────────────────────────────────────────────────────────────
\bibliographystyle{plain}
\begin{thebibliography}{99}

\bibitem{koza1992}
J.~R. Koza.
\newblock \emph{Genetic Programming: On the Programming of Computers by Means of Natural Selection}.
\newblock MIT Press, 1992.

\bibitem{udrescu2020}
S.-M. Udrescu and M.~Tegmark.
\newblock {AI Feynman}: A physics-inspired method for symbolic regression.
\newblock \emph{Science Advances}, 6(16):eaay2631, 2020.

\bibitem{cranmer2023}
M.~Cranmer.
\newblock Interpretable machine learning for science with {PySR} and {SymbolicRegression.jl}.
\newblock \emph{arXiv:2305.01582}, 2023.

\bibitem{brunton2016}
S.~L. Brunton, J.~L. Proctor, and J.~N. Kutz.
\newblock Discovering governing equations from data by sparse identification of nonlinear
  dynamical systems.
\newblock \emph{Proceedings of the National Academy of Sciences}, 113(15):3932--3937, 2016.

\bibitem{gompertz1825}
B.~Gompertz.
\newblock On the nature of the function expressive of the law of human mortality.
\newblock \emph{Philosophical Transactions of the Royal Society}, 115:513--583, 1825.

\bibitem{verhulst1838}
P.-F. Verhulst.
\newblock Notice sur la loi que la population poursuit dans son accroissement.
\newblock \emph{Correspondance Math\'ematique et Physique}, 10:113--121, 1838.

\bibitem{baranyi1994}
J.~Baranyi and T.~A. Roberts.
\newblock A dynamic approach to predicting bacterial growth in food.
\newblock \emph{International Journal of Food Microbiology}, 23(3--4):277--294, 1994.

\bibitem{richards1959}
F.~J. Richards.
\newblock A flexible growth function for empirical use.
\newblock \emph{Journal of Experimental Botany}, 10(2):290--301, 1959.

\bibitem{lindeberg1994}
T.~Lindeberg.
\newblock Scale-space theory: A basic tool for analysing structures at different scales.
\newblock \emph{Journal of Applied Statistics}, 21(1--2):225--270, 1994.

\bibitem{messenger2021}
D.~A. Messenger and D.~M. Bortz.
\newblock Weak {SINDy} for partial differential equations.
\newblock \emph{Journal of Computational Physics}, 443:110525, 2021.

\end{thebibliography}

\end{document}
